\documentclass[10pt,a4paper]{amsart}
\usepackage{amsmath,amssymb,amsthm,graphicx}
% Retained preamble for the Stacks Project chapter {CH} (shared project preamble).
% The chapter file pinned-original-source.tex only carries \input{{preamble}} on its first line,
% so no preamble text lives inside the chapter itself.  The Stacks Project uses one single shared
% preamble for every chapter: the snapshot stages it three times (sources/stacks-chow/preamble.tex,
% sources/stacks-etale-cohomology/preamble.tex and focused-seeds/AGP/sources/stacks-snapshot-20261011/
% preamble.tex) and those copies are identical in size.  The block below is copied mechanically,
% line for line, from sources/stacks-chow/preamble.tex lines 14-18 and 154-217: the Xy-pic setup and
% the native theorem declarations together with the mathematical macros.  Nothing was added, removed
% or redefined; the wrapper class and the amsmath/amssymb/amsthm/graphicx base packages are supplied
% by the assembler.
% For commutative diagrams we use Xy-pic
% Chapter: stacks-spaces-more-morphisms.tex -- More on Morphisms of Spaces.
\usepackage[all]{xy}

% We use 2cell for 2-commutative diagrams.
\xyoption{2cell}

% Theorem environments.
%
\theoremstyle{plain}
\newtheorem{theorem}[subsection]{Theorem}
\newtheorem{proposition}[subsection]{Proposition}
\newtheorem{lemma}[subsection]{Lemma}

\theoremstyle{definition}
\newtheorem{definition}[subsection]{Definition}
\newtheorem{example}[subsection]{Example}
\newtheorem{exercise}[subsection]{Exercise}
\newtheorem{situation}[subsection]{Situation}

\theoremstyle{remark}
\newtheorem{remark}[subsection]{Remark}
\newtheorem{remarks}[subsection]{Remarks}

\numberwithin{equation}{subsection}

% Macros
%
\def\lim{\mathop{\mathrm{lim}}\nolimits}
\def\colim{\mathop{\mathrm{colim}}\nolimits}
\def\Spec{\mathop{\mathrm{Spec}}}
\def\Hom{\mathop{\mathrm{Hom}}\nolimits}
\def\Ext{\mathop{\mathrm{Ext}}\nolimits}
\def\SheafHom{\mathop{\mathcal{H}\!\mathit{om}}\nolimits}
\def\SheafExt{\mathop{\mathcal{E}\!\mathit{xt}}\nolimits}
\def\Sch{\mathit{Sch}}
\def\Mor{\mathop{\mathrm{Mor}}\nolimits}
\def\Ob{\mathop{\mathrm{Ob}}\nolimits}
\def\Sh{\mathop{\mathit{Sh}}\nolimits}
\def\NL{\mathop{N\!L}\nolimits}
\def\CH{\mathop{\mathrm{CH}}\nolimits}
\def\proetale{{pro\text{-}\acute{e}tale}}
\def\etale{{\acute{e}tale}}
\def\QCoh{\mathit{QCoh}}
\def\Ker{\mathop{\mathrm{Ker}}}
\def\Im{\mathop{\mathrm{Im}}}
\def\Coker{\mathop{\mathrm{Coker}}}
\def\Coim{\mathop{\mathrm{Coim}}}

% Boxtimes
%
\DeclareMathSymbol{\boxtimes}{\mathbin}{AMSa}{"02}

%
% Macros for moduli stacks/spaces
%
\def\QCohstack{\mathcal{QC}\!\mathit{oh}}
\def\Cohstack{\mathcal{C}\!\mathit{oh}}
\def\Spacesstack{\mathcal{S}\!\mathit{paces}}
\def\Quotfunctor{\mathrm{Quot}}
\def\Hilbfunctor{\mathrm{Hilb}}
\def\Curvesstack{\mathcal{C}\!\mathit{urves}}
\def\Polarizedstack{\mathcal{P}\!\mathit{olarized}}
\def\Complexesstack{\mathcal{C}\!\mathit{omplexes}}
% \Pic is the operator that assigns to X its picard group, usage \Pic(X)
% \Picardstack_{X/B} denotes the Picard stack of X over B
% \Picardfunctor_{X/B} denotes the Picard functor of X over B
\def\Pic{\mathop{\mathrm{Pic}}\nolimits}
\def\Picardstack{\mathcal{P}\!\mathit{ic}}
\def\Picardfunctor{\mathrm{Pic}}
\def\Deformationcategory{\mathcal{D}\!\mathit{ef}}
\newenvironment{reference}{}{}
\newenvironment{slogan}{}{}
\newenvironment{history}{}{}
\usepackage{cleveref}
\crefname{originallemma}{original lemma}{original lemmas}
\Crefname{originallemma}{Original lemma}{Original lemmas}
\makeatletter
\newlabel{lemma-quasi-finite-separated-quasi-affine}{{02LR}{source}{}{sourcepointer-1}{}}
\newlabel{lemma-quasi-finite-separated-quasi-affine@cref}{{[originallemma][0][]02LR}{[1][0][]source}{}{}{}}
\makeatother
\crefname{originallemma}{original lemma}{original lemmas}
\Crefname{originallemma}{Original lemma}{Original lemmas}
\makeatletter
\newlabel{spaces-limits-lemma-descend-opens}{{0827}{source}{}{sourcepointer-2}{}}
\newlabel{spaces-limits-lemma-descend-opens@cref}{{[originallemma][0][]0827}{[1][0][]source}{}{}{}}
\makeatother
\crefname{originallemma}{original lemma}{original lemmas}
\Crefname{originallemma}{Original lemma}{Original lemmas}
\makeatletter
\newlabel{spaces-limits-lemma-directed-inverse-system-has-limit}{{07SF}{source}{}{sourcepointer-3}{}}
\newlabel{spaces-limits-lemma-directed-inverse-system-has-limit@cref}{{[originallemma][0][]07SF}{[1][0][]source}{}{}{}}
\makeatother
\crefname{originallemma}{original lemma}{original lemmas}
\Crefname{originallemma}{Original lemma}{Original lemmas}
\makeatletter
\newlabel{spaces-limits-lemma-finite-type-eventually-closed}{{0828}{source}{}{sourcepointer-4}{}}
\newlabel{spaces-limits-lemma-finite-type-eventually-closed@cref}{{[originallemma][0][]0828}{[1][0][]source}{}{}{}}
\makeatother
\crefname{originallemma}{original lemma}{original lemmas}
\Crefname{originallemma}{Original lemma}{Original lemmas}
\makeatletter
\newlabel{spaces-limits-lemma-integral-algebra-directed-colimit-finite}{{082C}{source}{}{sourcepointer-5}{}}
\newlabel{spaces-limits-lemma-integral-algebra-directed-colimit-finite@cref}{{[originallemma][0][]082C}{[1][0][]source}{}{}{}}
\makeatother
\crefname{originallemma}{original lemma}{original lemmas}
\Crefname{originallemma}{Original lemma}{Original lemmas}
\makeatletter
\newlabel{spaces-morphisms-lemma-factor-the-other-way}{{081U}{source}{}{sourcepointer-6}{}}
\newlabel{spaces-morphisms-lemma-factor-the-other-way@cref}{{[originallemma][0][]081U}{[1][0][]source}{}{}{}}
\makeatother
\title{Zariski's Main Theorem for spaces: quasi-finite separated maps factor through a finite space over qcqs targets}
\author{The Stacks Project Authors}
\date{stacks-project:spaces-more-morphisms.tex:local-snapshot-20261011; GNU FDL 1.2 or later}
\begin{document}
\maketitle
\noindent\textit{Editorial note.} Statement span is lines 7754-7768 of sources/stacks-spaces-more-morphisms/pinned-original-source.tex; proof span is lines 7770-7803, the complete author proof (32 body lines).\par\medskip
\noindent\textit{Editorial note.} Navigation stubs: 6 cited-but-not-included labels are supplied as non-load-bearing pointers with their pinned tags where the snapshot tag index has them.\par\medskip
\noindent\textit{Editorial note.} extra\_packages is empty on purpose: the retained preamble already loads xy with options and a second declaration clashes in the assembler wrapper.\par\medskip
\noindent\textit{Editorial note.} Text is retained verbatim, including the author's own typos and placeholders.\par\medskip
\noindent\textit{Editorial note.} Counted B side (both complete). The spaces ZMT card the derived-Chow and Stein machinery already lean on.\par\medskip
\noindent\textit{Allowed standard prerequisites (editorial).}\begin{itemize}
\item Explicit prerequisite boundary. Chapter lemma quasi-finite-separated-quasi-affine; Limits of Spaces: integral-algebra-directed-colimit-finite, descend-opens, directed-inverse-system-has-limit, finite-type-eventually-closed; Morphisms of Spaces, lemma factor-the-other-way.
\end{itemize}\medskip
\ifdefined\hypertarget\hypertarget{sourcepointer-1}{}\fi
\noindent\textit{Source reference.} Original lemma 02LR: sources/stacks-spaces-more-morphisms/pinned-original-source.tex, source lines 7724--7738. Cited by the retained statement or proof but not part of the collected spans; supplied only as a navigation pointer. The argument does not depend on its content.\par\medskip
\ifdefined\hypertarget\hypertarget{sourcepointer-2}{}\fi
\noindent\textit{Source reference.} Original lemma 0827: sources/stacks-spaces-more-morphisms/pinned-original-source.tex, source lines 7803--7803. Cited by the retained statement or proof but not part of the collected spans; supplied only as a navigation pointer. The argument does not depend on its content.\par\medskip
\ifdefined\hypertarget\hypertarget{sourcepointer-3}{}\fi
\noindent\textit{Source reference.} Original lemma 07SF: sources/stacks-spaces-more-morphisms/pinned-original-source.tex, source lines 7803--7803. Cited by the retained statement or proof but not part of the collected spans; supplied only as a navigation pointer. The argument does not depend on its content.\par\medskip
\ifdefined\hypertarget\hypertarget{sourcepointer-4}{}\fi
\noindent\textit{Source reference.} Original lemma 0828: sources/stacks-spaces-more-morphisms/pinned-original-source.tex, source lines 7803--7803. Cited by the retained statement or proof but not part of the collected spans; supplied only as a navigation pointer. The argument does not depend on its content.\par\medskip
\ifdefined\hypertarget\hypertarget{sourcepointer-5}{}\fi
\noindent\textit{Source reference.} Original lemma 082C: sources/stacks-spaces-more-morphisms/pinned-original-source.tex, source lines 7803--7803. Cited by the retained statement or proof but not part of the collected spans; supplied only as a navigation pointer. The argument does not depend on its content.\par\medskip
\ifdefined\hypertarget\hypertarget{sourcepointer-6}{}\fi
\noindent\textit{Source reference.} Original lemma 081U: sources/stacks-spaces-more-morphisms/pinned-original-source.tex, source lines 7803--7803. Cited by the retained statement or proof but not part of the collected spans; supplied only as a navigation pointer. The argument does not depend on its content.\par\medskip
\begin{lemma}[Zariski's Main Theorem]
\label{lemma-quasi-finite-separated-pass-through-finite}
Let $S$ be a scheme.
Let $f : X \to Y$ be a morphism of algebraic spaces over $S$.
Assume $f$ is quasi-finite and separated and assume that
$Y$ is quasi-compact and quasi-separated. Then there exists
a factorization
$$
\xymatrix{
X \ar[rd]_f \ar[rr]_j & & T \ar[ld]^\pi \\
& Y &
}
$$
where $j$ is a quasi-compact open immersion and $\pi$ is finite.
\end{lemma}

\begin{proof}
Let $X \to Y' \to Y$ be as in the conclusion of
Lemma \ref{lemma-quasi-finite-separated-quasi-affine}.
By
Limits of Spaces, Lemma
\ref{spaces-limits-lemma-integral-algebra-directed-colimit-finite}
we can write
$\nu_*\mathcal{O}_{Y'} = \colim_{i \in I} \mathcal{A}_i$ as a
directed colimit of finite quasi-coherent $\mathcal{O}_X$-algebras
$\mathcal{A}_i \subset \nu_*\mathcal{O}_{Y'}$. Then
$\pi_i : T_i = \underline{\Spec}_Y(\mathcal{A}_i) \to Y$
is a finite morphism for each $i$.
Note that the transition morphisms $T_{i'} \to T_i$ are affine
and that $Y' = \lim T_i$.

\medskip\noindent
By Limits of Spaces, Lemma \ref{spaces-limits-lemma-descend-opens}
there exists an $i$ and a quasi-compact open
$U_i \subset T_i$ whose inverse image in $Y'$ equals
$f'(X)$. For $i' \geq i$ let $U_{i'}$ be the inverse image
of $U_i$ in $T_{i'}$. Then $X \cong f'(X) = \lim_{i' \geq i} U_{i'}$, see
Limits of Spaces, Lemma
\ref{spaces-limits-lemma-directed-inverse-system-has-limit}.
By
Limits of Spaces, Lemma
\ref{spaces-limits-lemma-finite-type-eventually-closed} we see that
$X \to U_{i'}$ is a closed immersion for some $i' \geq i$.
(In fact $X \cong U_{i'}$ for sufficiently
large $i'$ but we don't need this.) Hence $X \to T_{i'}$ is an immersion. By
Morphisms of Spaces, Lemma
\ref{spaces-morphisms-lemma-factor-the-other-way}
we can factor this as $X \to T \to T_{i'}$ where the first arrow
is an open immersion and the second a closed immersion. Thus we win.
\end{proof}

\end{document}
